\documentclass[12pt, a4paper]{ctexart}
\usepackage{fontspec}
\usepackage{xcolor}
\usepackage{hyperref}
\usepackage{geometry}
\usepackage{enumitem}
\usepackage{parskip}
\usepackage{titlesec}
\usepackage{tcolorbox}
\tcbuselibrary{breakable}
\usepackage{setspace}
\usepackage{listings}

\geometry{left=2.5cm, right=2.5cm, top=2.5cm, bottom=2.5cm}

\setmainfont{Times New Roman}
\setCJKmainfont{SimSun}

\hypersetup{
    colorlinks=true,
    linkcolor=blue!70!black,
    urlcolor=cyan,
    pdftitle={Java集合-逐题精讲解义},
    pdfauthor={专升本-fifine老师}
}

\definecolor{myblue}{RGB}{30,100,200}
\definecolor{lightblue}{RGB}{230,240,255}
\definecolor{darkblue}{RGB}{0,50,120}

\newtcolorbox{bluebox}[1][]{
    colback=lightblue,
    colframe=myblue,
    coltitle=darkblue,
    fonttitle=\bfseries,
    title={#1},
    boxrule=0.8pt,
    arc=4pt,
    left=6pt,
    right=6pt,
    top=6pt,
    bottom=6pt,
    before skip=8pt,
    after skip=8pt,
    breakable
}

\usepackage{etoolbox}
\BeforeBeginEnvironment{bluebox}{\par\vfill}%
\AfterEndEnvironment{bluebox}{\par\vfill}%

\titleformat{\section}
  {\normalfont\Large\bfseries\color{myblue}}{\thesection}{1em}{}
\titlespacing*{\section}{0pt}{3.5ex plus 1ex minus .2ex}{1.5ex plus .2ex}

\title{\textcolor{myblue}{\textbf{Java集合 - 逐题精讲解义}}}
\author{专升本-fifine老师 教学组}
\date{2026年03月02日}

\lstset{
    language=Java,
    basicstyle=\ttfamily\small,
    keywordstyle=\color{blue},
    commentstyle=\color{green!60!black},
    stringstyle=\color{orange},
    numbers=left,
    numberstyle=\tiny\color{gray},
    frame=single,
    breaklines=true,
    columns=fullflexible,
    showstringspaces=false
}

\newcommand{\code}[1]{\texttt{#1}}

\begin{document}

\maketitle
\thispagestyle{empty}
\newpage

\section{集合}

\begin{enumerate}[label=\textbf{\arabic*.}]

	\item \textbf{以下关于Java集合框架的说法正确的是}
	      \begin{enumerate}[label=\Alph*.]
		      \item List允许重复元素
		      \item Set允许重复元素
		      \item Map存储键值对，键可以重复
		      \item 以上都不对
	      \end{enumerate}

	      \begin{bluebox}{答案与深度解析}
		      \textbf{答案：A. List允许重复元素}

		      \textbf{【核心认知框架】}

		      \textbf{【选项原理深度拆解】}
					  \begin{itemize}[label={}, leftmargin=*, nosep]

		      
		      \item \textbf{A. List允许重复元素 — 正确}
					  \begin{itemize}[label={}, nosep]
		      
			      \item List是有序集合，允许重复元素
			      \item 典型实现：ArrayList、LinkedList
		      

					  \end{itemize}
		      \item \textbf{B. Set允许重复元素 — 错误}
					  \begin{itemize}[label={}, nosep]
		      
			      \item Set是不允许重复的集合
			      \item 判断重复：equals()和hashCode()
		      

					  \end{itemize}
		      \item \textbf{C. Map键可以重复 — 错误}
					  \begin{itemize}[label={}, nosep]
		      
			      \item Map的键(key)必须唯一
			      \item 值(value)可以重复
		      

					  \end{itemize}
		      \item \textbf{D. 以上都不对 — 错误}
					  \begin{itemize}[label={}, nosep]
		      
			      \item A选项正确
		      
		      

		      
					  \end{itemize}
					  \end{itemize}
\textbf{【应背诵知识点】}

		      集合特性速查：
		      \begin{tabular}{|c|c|c|c|}
			      \hline
			      接口   & 有序 & 允许重复 & 典型实现                  \\
			      \hline
			      List & ✓  & ✓    & ArrayList, LinkedList \\
			      \hline
			      Set  & -  & ✗    & HashSet, TreeSet      \\
			      \hline
			      Map  & -  & 键✗值✓ & HashMap, TreeMap      \\
			      \hline
		      \end{tabular}
	      \end{bluebox}

	      \vspace{1em}

	\item \textbf{在Java中，以下哪个集合实现了链表结构？}
	      \begin{enumerate}[label=\Alph*.]
		      \item ArrayList
		      \item LinkedList
		      \item HashSet
		      \item HashMap
	      \end{enumerate}

	      \begin{bluebox}{答案与深度解析}
		      \textbf{答案：B. LinkedList}

		      \textbf{【核心认知框架】}

		      \textbf{【选项原理深度拆解】}
					  \begin{itemize}[label={}, leftmargin=*, nosep]

		      
		      \item \textbf{A. ArrayList — 错误}
					  \begin{itemize}[label={}, nosep]
		      
			      \item 基于动态数组实现
			      \item 随机访问效率高，插入删除效率低
		      

					  \end{itemize}
		      \item \textbf{B. LinkedList — 正确}
					  \begin{itemize}[label={}, nosep]
		      
			      \item 基于双向链表实现
			      \item 插入删除效率高，随机访问效率低
		      

					  \end{itemize}
		      \item \textbf{C. HashSet — 错误}
					  \begin{itemize}[label={}, nosep]
		      
			      \item 基于哈希表实现
		      

					  \end{itemize}
		      \item \textbf{D. HashMap — 错误}
					  \begin{itemize}[label={}, nosep]
		      
			      \item 基于哈希表+链表/红黑树实现
		      
		      

		      
					  \end{itemize}
					  \end{itemize}
\textbf{【应背诵知识点】}

		      数据结构对应：
		      \begin{itemize}[label={}, nosep]
			      \item 数组：ArrayList
			      \item 链表：LinkedList
			      \item 哈希表：HashSet、HashMap
			      \item 树：TreeSet、TreeMap
			      \item 链表+哈希表：LinkedHashSet、LinkedHashMap
		      \end{itemize}
	      \end{bluebox}

	      \vspace{1em}

	\item \textbf{以下哪个集合是无序且不允许重复元素的？}
	      \begin{enumerate}[label=\Alph*.]
		      \item ArrayList
		      \item HashSet
		      \item LinkedHashSet
		      \item TreeSet
	      \end{enumerate}

	      \begin{bluebox}{答案与深度解析}
		      \textbf{答案：B. HashSet}

		      \textbf{【核心认知框架】}

		      \textbf{【选项原理深度拆解】}
					  \begin{itemize}[label={}, leftmargin=*, nosep]

		      
		      \item \textbf{A. ArrayList — 错误}
					  \begin{itemize}[label={}, nosep]
		      
			      \item 有序，允许重复
		      

					  \end{itemize}
		      \item \textbf{B. HashSet — 正确}
					  \begin{itemize}[label={}, nosep]
		      
			      \item 无序（不保证顺序）
			      \item 不允许重复（基于hashCode和equals）
		      

					  \end{itemize}
		      \item \textbf{C. LinkedHashSet — 错误}
					  \begin{itemize}[label={}, nosep]
		      
			      \item 有序（保持插入顺序）
		      

					  \end{itemize}
		      \item \textbf{D. TreeSet — 错误}
					  \begin{itemize}[label={}, nosep]
		      
			      \item 有序（基于红黑树，排序顺序）
		      
		      

		      
					  \end{itemize}
					  \end{itemize}
\textbf{【应背诵知识点】}

		      Set实现类对比：
		      \begin{itemize}[label={}, nosep]
			      \item HashSet：无序，性能最好
			      \item LinkedHashSet：插入顺序
			      \item TreeSet：自然顺序/自定义排序
		      \end{itemize}
	      \end{bluebox}

	      \vspace{1em}

	\item \textbf{以下关于HashMap的说法错误的是}
	      \begin{enumerate}[label=\Alph*.]
		      \item 存储键值对
		      \item 键不允许重复
		      \item 是线程安全的
		      \item 基于哈希表实现
	      \end{enumerate}

	      \begin{bluebox}{答案与深度解析}
		      \textbf{答案：C. 是线程安全的}

		      \textbf{【核心认知框架】}

		      \textbf{【选项原理深度拆解】}
					  \begin{itemize}[label={}, leftmargin=*, nosep]

		      
		      \item \textbf{A. 存储键值对 — 正确}
					  \begin{itemize}[label={}, nosep]
		      
			      \item Map接口的核心特性
		      

					  \end{itemize}
		      \item \textbf{B. 键不允许重复 — 正确}
					  \begin{itemize}[label={}, nosep]
		      
			      \item 键唯一，值可重复
		      

					  \end{itemize}
		      \item \textbf{C. 线程安全 — 错误}
					  \begin{itemize}[label={}, nosep]
		      
			      \item HashMap非线程安全
			      \item Hashtable和ConcurrentHashMap是线程安全的
		      

					  \end{itemize}
		      \item \textbf{D. 基于哈希表实现 — 正确}
					  \begin{itemize}[label={}, nosep]
		      
			      \item JDK8+：数组+链表+红黑树
		      
		      

		      
					  \end{itemize}
					  \end{itemize}
\textbf{【应背诵知识点】}

		      HashMap vs Hashtable：
		      \begin{itemize}[label={}, nosep]
			      \item 线程安全：Hashtable(✓) vs HashMap(✗)
			      \item null支持：HashMap允许null键值，Hashtable不允许
			      \item 性能：HashMap更好
			      \item 迭代器：HashMap用Iterator，Hashtable用Enumeration
		      \end{itemize}
	      \end{bluebox}

	      \vspace{1em}

	\item \textbf{以下哪个方法可以判断一个集合是否为空？}
	      \begin{enumerate}[label=\Alph*.]
		      \item isEmpty()
		      \item isBlank()
		      \item isNullOrEmpty()
		      \item isNotNull()
	      \end{enumerate}

	      \begin{bluebox}{答案与深度解析}
		      \textbf{答案：A. isEmpty()}

		      \textbf{【核心认知框架】}

		      \textbf{【选项原理深度拆解】}
					  \begin{itemize}[label={}, leftmargin=*, nosep]

		      
		      \item \textbf{A. isEmpty() — 正确}
					  \begin{itemize}[label={}, nosep]
		      
			      \item Collection接口定义的方法
			      \item 所有集合都支持
		      

					  \end{itemize}
		      \item \textbf{B. isBlank() — 错误}
					  \begin{itemize}[label={}, nosep]
		      
			      \item String类的方法（判断空白字符串）
		      

					  \end{itemize}
		      \item \textbf{C. isNullOrEmpty() — 错误}
					  \begin{itemize}[label={}, nosep]
		      
			      \item 不是集合的标准方法
		      

					  \end{itemize}
		      \item \textbf{D. isNotNull() — 错误}
					  \begin{itemize}[label={}, nosep]
		      
			      \item 不是集合的方法
		      
		      
	      
	      
					  \end{itemize}
					  \end{itemize}
\end{bluebox}

	      \vspace{1em}

	\item \textbf{以下关于Java中Lambda表达式的说法正确的是}
	      \begin{enumerate}[label=\Alph*.]
		      \item 可以简化匿名内部类的写法
		      \item 不支持函数式接口
		      \item 不能作为方法参数
		      \item 性能比普通方法差
	      \end{enumerate}

	      \begin{bluebox}{答案与深度解析}
		      \textbf{答案：A. 可以简化匿名内部类的写法}

		      \textbf{【核心认知框架】}

		      \textbf{【选项原理深度拆解】}
					  \begin{itemize}[label={}, leftmargin=*, nosep]

		      
		      \item \textbf{A. 简化匿名内部类 — 正确}
					  \begin{itemize}[label={}, nosep]
		      
			      \item Lambda是函数式接口的简洁实现
			      \item (参数) -> 表达式/语句块
		      

					  \end{itemize}
		      \item \textbf{B. 不支持函数式接口 — 错误}
					  \begin{itemize}[label={}, nosep]
		      
			      \item Lambda就是用于实现函数式接口的
		      

					  \end{itemize}
		      \item \textbf{C. 不能作为方法参数 — 错误}
					  \begin{itemize}[label={}, nosep]
		      
			      \item Lambda可以作为参数传递
		      

					  \end{itemize}
		      \item \textbf{D. 性能比普通方法差 — 错误}
					  \begin{itemize}[label={}, nosep]
		      
			      \item 在运行时没有明显性能差异
			      \item JIT编译后会内联
		      
		      

		      
					  \end{itemize}
					  \end{itemize}
\textbf{【应背诵知识点】}

		      Lambda表达式：
		      \begin{itemize}[label={}, nosep]
			      \item 语法：(parameters) -> expression 或 (parameters) -> \{ statements; \}
			      \item 函数式接口：只包含一个抽象方法的接口
			      \item 常见函数式接口：Runnable、Callable、Comparator、Function
			      \item 方法引用：ClassName::method
		      \end{itemize}
	      \end{bluebox}

	      \vspace{1em}

	\item \textbf{以下哪个是函数式接口？}
	      \begin{enumerate}[label=\Alph*.]
		      \item 接口中有多个抽象方法
		      \item 接口中有一个抽象方法
		      \item 接口中有静态方法
		      \item 接口中有默认方法
	      \end{enumerate}

	      \begin{bluebox}{答案与深度解析}
		      \textbf{答案：B. 接口中有一个抽象方法}

		      \textbf{【核心认知框架】}

		      \textbf{【选项原理深度拆解】}
					  \begin{itemize}[label={}, leftmargin=*, nosep]

		      
		      \item \textbf{A. 多个抽象方法 — 错误}
					  \begin{itemize}[label={}, nosep]
		      
			      \item 这不是函数式接口
		      

					  \end{itemize}
		      \item \textbf{B. 一个抽象方法 — 正确}
					  \begin{itemize}[label={}, nosep]
		      
			      \item 函数式接口定义
			      \item 可以包含default和static方法
		      

					  \end{itemize}
		      \item \textbf{C/D. 静态/默认方法 — 错误}
					  \begin{itemize}[label={}, nosep]
		      
			      \item 这些不影响函数式接口的判定
		      
		      

		      
					  \end{itemize}
					  \end{itemize}
\textbf{【应背诵知识点】}

		      函数式接口：
		      \begin{itemize}[label={}, nosep]
			      \item @FunctionalInterface注解（可选）
			      \item 只能有一个抽象方法
			      \item 可以有多个default和static方法
			      \item 常见：Runnable、Supplier<T>、Consumer<T>、Function<T,R>、Predicate<T>
		      \end{itemize}
	      \end{bluebox}

	      \vspace{1em}

	\item \textbf{以下关于Java中方法引用的说法错误的是}
	      \begin{enumerate}[label=\Alph*.]
		      \item 可以减少代码量
		      \item 有四种类型的方法引用
		      \item 不能引用静态方法
		      \item 是Lambda表达式的一种简洁写法
	      \end{enumerate}

	      \begin{bluebox}{答案与深度解析}
		      \textbf{答案：C. 不能引用静态方法}

		      \textbf{【核心认知框架】}

		      \textbf{【选项原理深度拆解】}
					  \begin{itemize}[label={}, leftmargin=*, nosep]

		      
		      \item \textbf{A. 减少代码量 — 正确}
					  \begin{itemize}[label={}, nosep]
		      
			      \item 方法引用是Lambda的简化形式
		      

					  \end{itemize}
		      \item \textbf{B. 四种类型 — 正确}
					  \begin{itemize}[label={}, nosep]
		      
			      \item 静态方法引用：ClassName::staticMethod
			      \item 实例方法引用（对象）：instance::method
			      \item 实例方法引用（类型）：ClassName::method
			      \item 构造方法引用：ClassName::new
		      

					  \end{itemize}
		      \item \textbf{C. 不能引用静态方法 — 错误}
					  \begin{itemize}[label={}, nosep]
		      
			      \item 可以引用！如System.out::println
		      

					  \end{itemize}
		      \item \textbf{D. Lambda简洁写法 — 正确}
					  \begin{itemize}[label={}, nosep]
		      
			      \item 方法引用本质上就是Lambda表达式
		      
		      
	      
	      
					  \end{itemize}
					  \end{itemize}
\end{bluebox}

	      \vspace{1em}

	\item \textbf{以下关于Java中Stream流的说法正确的是}
	      \begin{enumerate}[label=\Alph*.]
		      \item 可以对集合进行高效的操作
		      \item 只能顺序处理数据
		      \item 不支持并行处理
		      \item 不能进行数据过滤
	      \end{enumerate}

	      \begin{bluebox}{答案与深度解析}
		      \textbf{答案：A. 可以对集合进行高效的操作}

		      \textbf{【核心认知框架】}

		      \textbf{【选项原理深度拆解】}
					  \begin{itemize}[label={}, leftmargin=*, nosep]

		      
		      \item \textbf{A. 高效操作集合 — 正确}
					  \begin{itemize}[label={}, nosep]
		      
			      \item Stream提供函数式操作
			      \item 支持链式调用
		      

					  \end{itemize}
		      \item \textbf{B. 只能顺序处理 — 错误}
					  \begin{itemize}[label={}, nosep]
		      
			      \item Stream支持并行处理：parallelStream()
		      

					  \end{itemize}
		      \item \textbf{C. 不支持并行处理 — 错误}
					  \begin{itemize}[label={}, nosep]
		      
			      \item 支持！parallelStream()利用多核
		      

					  \end{itemize}
		      \item \textbf{D. 不能过滤 — 错误}
					  \begin{itemize}[label={}, nosep]
		      
			      \item filter()用于过滤数据
		      
		      

		      
					  \end{itemize}
					  \end{itemize}
\textbf{【应背诵知识点】}

		      Stream操作分类：
		      \begin{itemize}[label={}, nosep]
			      \item 中间操作：filter、map、sorted、distinct...
			      \item 终端操作：forEach、collect、reduce、count...
			      \item 惰性求值：中间操作不执行，等终端操作触发
			      \item parallelStream()：并行流，提升性能
		      \end{itemize}
	      \end{bluebox}

	      \vspace{1em}

	\item \textbf{以下哪个方法用于创建Stream流？}
	      \begin{enumerate}[label=\Alph*.]
		      \item stream()
		      \item createStream()
		      \item makeStream()
		      \item generateStream()
	      \end{enumerate}

	      \begin{bluebox}{答案与深度解析}
		      \textbf{答案：A. stream()}

		      \textbf{【核心认知框架】}

		      \textbf{【选项原理深度拆解】}
					  \begin{itemize}[label={}, leftmargin=*, nosep]

		      
		      \item \textbf{A. stream() — 正确}
					  \begin{itemize}[label={}, nosep]
		      
			      \item Collection接口的默认方法
			      \item collection.stream()
		      

					  \end{itemize}
		      \item \textbf{B/C/D — 错误}
					  \begin{itemize}[label={}, nosep]
		      
			      \item 这些方法不存在
		      
		      

		      
					  \end{itemize}
					  \end{itemize}
\textbf{【应背诵知识点】}

		      创建Stream：
		      \begin{itemize}[label={}, nosep]
			      \item 集合：collection.stream() / collection.parallelStream()
			      \item 数组：Arrays.stream(array)
			      \item 静态方法：Stream.of()、Stream.iterate()、Stream.generate()
			      \item 文件：Files.lines()
		      \end{itemize}
	      \end{bluebox}

	      \vspace{1em}

	\item \textbf{以下关于Java中Stream流的终端操作的说法错误的是}
	      \begin{enumerate}[label=\Alph*.]
		      \item 终端操作会触发流的执行
		      \item forEach()是终端操作
		      \item 终端操作只能执行一次
		      \item 终端操作后可以继续进行中间操作
	      \end{enumerate}

	      \begin{bluebox}{答案与深度解析}
		      \textbf{答案：D. 终端操作后可以继续进行中间操作}

		      \textbf{【核心认知框架】}

		      \textbf{【选项原理深度拆解】}
					  \begin{itemize}[label={}, leftmargin=*, nosep]

		      
		      \item \textbf{A. 触发执行 — 正确}
					  \begin{itemize}[label={}, nosep]
		      
			      \item 终端操作触发中间操作执行
		      

					  \end{itemize}
		      \item \textbf{B. forEach是终端操作 — 正确}
					  \begin{itemize}[label={}, nosep]
		      
			      \item forEach是常见终端操作
		      

					  \end{itemize}
		      \item \textbf{C. 只能执行一次 — 正确}
					  \begin{itemize}[label={}, nosep]
		      
			      \item 流是消耗性的，一次使用后失效
		      

					  \end{itemize}
		      \item \textbf{D. 可继续中间操作 — 错误}
					  \begin{itemize}[label={}, nosep]
		      
			      \item 终端操作后流被关闭，不能再操作
		      
		      

		      
					  \end{itemize}
					  \end{itemize}
\textbf{【应背诵知识点】}

		      Stream特点：
		      \begin{itemize}[label={}, nosep]
			      \item 惰性求值：中间操作不立即执行
			      \item 一次性：终端操作后流关闭
			      \item 内部迭代：不同于传统的外部迭代
			      \item 常见终端操作：forEach、collect、reduce、count、min、max
		      \end{itemize}
	      \end{bluebox}

\end{enumerate}

\end{document}
